% -----------------------------------------------------------
\section{A short introduction to Climbing the Rainbows}
% -----------------------------------------------------------

I ran through a bunch of different beginnings for this piece, my first ever, but didn’t find anything I liked. So let me give this introduction:

Welcome to my essay newsletter, Climbing the Rainbows. Climbing the Rainbows is an essay series about the restored gospel, exploring it and its connections to science, relationships, professional and creative work, literature, and other aspects of life.

End of short introduction.

Now for the longer one.

% -----------------------------------------------------------
\section{Climbing the Rainbows: not just science!}
% -----------------------------------------------------------

In the next section I’m going to explain why I’m writing these essays and where all this comes from. I’m going to say a bunch of things about science, including some complaints. But I want you to know that these essays won’t just be about science—they’ll be about a lot of other things too. It just looks like I only care about the relationship between science and religion because the origin for this newsletter starts with that relationship.

% -----------------------------------------------------------
\section{So why write this stuff? Like really, why?}
% -----------------------------------------------------------

Yeah I ask myself that like all the time. 

I’ve been meaning to start these essays for a while, actually. For various reasons I just kept putting them off: “I haven’t taken enough notes yet.” “I need to exercise.” “I have to explain to my kids again why the new \textit{Star Wars} movies suck so bad.” But a while ago I had an experience that pushed me to get serious about working through these ideas and writing them down.

A friend messaged me, saying he’d been having a discussion with some other friends of his. They’d been talking about the gospel and one of them was explaining why he’d left the LDS church. Although at one time this person had spiritual experiences, he’d come to believe that those experiences were just illusions—they came not from God but from his brain’s own attempt to make sense of his membership in a certain group. The experiences therefore didn’t represent anything real and s o couldn’t justify his previous religious beliefs. This explanation also included a scientific-looking discussion of how “religious” experiences are valuable from an evolutionary perspective, but again, they weren’t \textit{real}; they didn’t mean God was communicating with him.

My friend asked for my take on this attempt to explain away one’s spiritual experiences. It wasn’t the first time I’ve come across this argument; I’ve seen other people make it before. I think it’s pretty interesting, and I see why so many people find it attractive.

I also think it’s wrong. It’s based on a misunderstanding of both neuroscience and evolutionary psychology.

The trouble is that the misunderstanding is really hard to see—so hard that many scientists don’t see it, including even some neuroscientists. Since the problem is very deep, the argument about spiritual experiences looks plausible on the surface and has a seductive appeal, especially for those who like to explain things in scientific terms.

Even so, it’s still wrong. So when my friend asked what I thought, I sat down to say why.

I wrote a few thousand words, both about why the problem is hard to see and why it makes the argument against spiritual experiences so ineffective. The things I said apply not only to the cognitive neuroscience of religion but to other areas of study as well; the problem undermines other scientific claims too. I wrote about the assumptions that frame the ways we think, and how these assumptions have a powerful effect on the conclusions we reach.

When I was done I sent the response to my friend, and that was that. Or at least, that’s what I thought at the time.

% -----------------------------------------------------------
\section{You haven’t said yet why you’re writing all this stuff}
% -----------------------------------------------------------

The more I thought about it, though, the more something bothered me about the whole experience. You’ve got someone who’s intelligent and educated—this particular person, the one explaining away his spiritual experiences, has a PhD. And that person is using the skills they gained in their professional training to make judgments about their religious experience. Those judgments then lead to life-changing decisions about how to live. But by their own ignorance they end up making those decisions for very poor reasons.

The least you can do is make those decisions for good reasons!

I’ve seen this pattern play out in the lives of many people, both with matters of faith and with other things too. They’re good people and they have good intentions. But because they don’t know some specific things about the way a certain science works, for example, or because they have false assumptions about the way the gospel operates, they’re in the dark about their real situation. They misunderstand the information available to them in evaluating their spiritual life and religious beliefs.

The worst part is that this ignorance is almost never their fault. They’re caught between two weak points, one coming from science and one from religion.

The first weak point is the public dissemination of science. Many scientific communities—psychology and neuroscience in particular—do a horrendous job of sharing information with non-experts about the true importance of their findings. And since fields like psychology and neuroscience are much more relevant to religious experience than, say, engineering, the people who receive the information know that it matters for their faith, but they have no idea how to evaluate it. They therefore can’t tell good arguments from bad ones. This is the first weak point.

The second weak point is on the religious side. In general, Christian religions in America do a horrendous job of helping their congregations to understand and answer problems with their beliefs. One thing these religions are unusually bad at is identifying the assumptions that people bring to their religious experience. But being bad at that is very, very bad, because those assumptions turn out to have incredible power to shape what you think about your faith.

As a result, many people end up in an unhappy place between these two weak points: scientists and other educated thought leaders are telling them one thing, and no one is around to help them examine those arguments; and on the other side, no one is helping them examine their own religious beliefs, either, to see which are worth keeping and which should head for the trash. It’s no wonder that so many people just end up claiming that everything is relative and that religion is just a matter of how you choose to live your life.

Now I’m not saying that I have an answer for every belief problem. I don’t. I’m not saying that I have a magic formula which will help people stay faithful to religious experience and remain an active participant in some church. There is no magic formula for that.[FOOTNOTE] But these essays will offer what I do have, which is tools to understand the insights that non-religious sources have for religious experience. I also have some perspectives on my own beliefs and doctrine that have helped me place them in the context of everything else I think I know about the world.

The purpose of Climbing the Rainbows is to share those tools and perspectives. That’s it. All I’ll do is offer them as possibilities. You can judge how they would change the way you see things and whether they might be beneficial.

% -----------------------------------------------------------
\section{More on the “not just science” part}
% -----------------------------------------------------------

A lot of those tools and perspectives don’t have anything to do with science, and so there’s tons more to say about other topics too. So, so many things can and should influence the way we think about the restored gospel of Jesus Christ. It connects to every other aspect of life. How could it not, after all? It claims to be the true religion, and if it is true, it better connect to everything else. It would be weird if it didn’t.

One of the things we’ll have most to say about is revelation: at the personal level, at the level of groups, at the level of the whole world. I talk about it so much because I’m convinced that it’s the most important topic in all of LDS thought. If you have questions about the gospel or its connections to other things, the best hope you have to answer them is to study revelation and the principles which govern its creation and reception. Later I’ll explain why I think that and will discuss those principles. Like a lot.

And there will be other stuff too. A lot about work and parenting, which are the two things that take most of my time. Also relationships and music and art. And video games. There will definitely be video games.

The point is that all these things are important parts of our lives. And since the gospel connects to them, we owe it to ourselves to treat these topics carefully and seriously.

% -----------------------------------------------------------
\section{End of long introduction}
% -----------------------------------------------------------

That’s it for now. This is Climbing the Rainbows, an essay series about the gospel and its connections to—everything, apparently.

Up next week: how not to explain away your religious experiences.